\section{A localized mass perturbation on the three-sphere}
\label{sec:mass}
In Newtonian (Poisson) gauge, take a scalar perturbation with negligible
anisotropic stress:
\begin{equation}\label{eq:perturbed}
 \dd s^2=A^2(\eta)\left[-(1+2\ph)\dd\eta^2
 +(1-2\ph)\gamma_{ij}\dd x^i\dd x^j\right]+\order(\epsilon^2).
\end{equation}
The potential $\ph=\Phi/c^2$ is dimensionless and of order $\epsilon$.
It perturbs both the lapse and the spatial metric. Curved-background scalar
perturbations are an established way to describe inhomogeneities
\citep{Canay2022}. We specify the source through the linear Einstein
constraints rather than inserting a flat-space $-GM/r$ potential globally.

\subsection{Evolution and reconstructed matter}
The linear dust equation for the potential, with unit positive spatial
curvature, is
\begin{equation}\label{eq:bardeen}
 \ph''+3\Hc\ph'+(2\Hc'+\Hc^2-1)\ph=0.
\end{equation}
The comoving density contrast $\Delta$, velocity potential $V$, and
Newtonian-gauge density contrast $\delta_N$ obey
\begin{align}
 \Delta&=\frac{2}{3(\Hc^2+1)}(\nabla_\gamma^2+3)\ph,
       \label{eq:constraint-density}\\
 V&=-\frac{2}{3(\Hc^2+1)}(\ph'+\Hc\ph),
       \label{eq:constraint-velocity}\\
 \delta_N&=\Delta+3\Hc V.\label{eq:constraint-gauge}
\end{align}
These are the first-order dust relations, specialized to the conventions
above \citep[Eq.~(14b)--(16)]{UgglaWainwright2013}. Our convention for the
dimensionless peculiar velocity is $v^i/c=\nabla_\gamma^i V$ at this order.
Using Eq.~\eqref{eq:hubble} directly verifies
\begin{equation}\label{eq:conservation}
 V'+\Hc V=-\ph,\qquad \delta_N'=-\nabla_\gamma^2V+3\ph'.
\end{equation}
The accompanying symbolic script checks both identities without selecting
one particular spatial eigenmode.

Write $\ph(\eta,\chi)=\epsilon g(\eta)F(\cos\chi)$, with $g(\eta_o)=1$.
Equation~\eqref{eq:bardeen} reduces for Eq.~\eqref{eq:dust} to
\begin{equation}\label{eq:growth}
 g''+3\cot(\eta/2)g'-2g=0.
\end{equation}
We select the regular growing solution by starting with a multiple of
\begin{equation}\label{eq:growth-series}
 1+\frac{\eta^2}{7}+\frac{\eta^4}{84}+\frac{\eta^6}{1320}
 +\order(\eta^8)
\end{equation}
at $\eta_i=10^{-3}$, and normalize at reception. An independent numerical
check uses the reduction-of-order expression
\begin{equation}
 g(\eta)\propto \frac{\Hc(\eta)}{A^2(\eta)}
 \int_0^\eta\frac{A(\tilde\eta)}{\Hc^2(\tilde\eta)}\,\dd\tilde\eta.
\end{equation}

\subsection{Smoothness, compensation, and the low spatial modes}
Let $\mu=\cos\chi$ and $\kappa=\sigma^{-2}$. Define
\begin{align}
 Q(\mu)&=\exp[\kappa(\mu-1)],\\
 q_0&=\langle Q\rangle=\frac{2e^{-\kappa}I_1(\kappa)}{\kappa},
 &q_1&=4\langle\mu Q\rangle=\frac{8e^{-\kappa}I_2(\kappa)}{\kappa},\\
 F(\mu)&=-\frac{Q(\mu)-q_0-q_1\mu}{1-q_0-q_1}.
 \label{eq:profile}
\end{align}
Here $I_n$ is the modified Bessel function and the normalized three-sphere
average of a radial function is
\begin{equation}
 \langle f\rangle=\frac{2}{\pi}\int_0^\pi f(\cos\chi)\sin^2\chi\,\dd\chi.
\end{equation}
Since $\langle\mu^2\rangle=1/4$, this construction gives
$F(1)=-1$, $\langle F\rangle=0$, and $\langle\mu F\rangle=0$.
It removes the homogeneous mode and the rotationally symmetric dipole
mode, avoiding a background redefinition and the exceptional lowest
scalar modes. Smoothness at both poles follows from smooth dependence
on $\cos\chi$.

The sphere Laplacian is computed without singular polar divisions:
\begin{equation}\label{eq:laplace}
 \nabla_\gamma^2F=(1-\mu^2)F_{,\mu\mu}-3\mu F_{,\mu}.
\end{equation}
For the specified parameters, Eq.~\eqref{eq:constraint-density} has a
positive central core and negative surrounding regions. Its integral
over the entire background three-sphere vanishes. The localized mass
is therefore an excess relative to the homogeneous dust distribution,
with compensation; it is not extra net mass added to an unchanged
compact background. The broad tail from the low-mode removal is included
in all calculations.

For a quantitative mass label, define the positive-core excess using
the comoving density contrast and background volume element:
\begin{equation}\label{eq:mass}
 M_+(\eta)=4\pi\bar\rho A^3
 \int_0^{\chi_c}\Delta(\eta,\chi)\sin^2\chi\,\dd\chi,
\end{equation}
where $\chi_c$ is the first zero of $\Delta$. This is a specified linear
background-relative mass diagnostic, not a nonlinear quasi-local mass.
Its evolving value is appropriate to growing dust structure.
